\section{Background}
\label{sec:background}

\subsection{RFC 8785: JSON Canonicalization Scheme}

\rfcjcs~\cite{rfc8785} defines a deterministic JSON serialization by
specifying: (1)~lexicographic key ordering using UTF-16 code unit
comparison, (2)~no insignificant whitespace, and (3)~floating-point
formatting matching \ecma \texttt{Number::toString}.

The floating-point requirement is the most technically demanding.  It
requires producing the \emph{shortest decimal representation} that
round-trips to the original IEEE~754 binary64 value under
roundTiesToEven, then applying the \ecma formatting policy (exponential
notation thresholds, sign handling, integer elision rules).

\subsection{ECMA-262 Number Formatting}

\ecma~\cite{ecma262} \S6.1.6.1.20 specifies \texttt{Number::toString} as
a two-stage process (formerly \S7.1.12.1 in older editions):

\begin{enumerate}
\item \textbf{Digit generation}: find the shortest decimal significand
      $d$ and exponent $n$ such that $d \times 10^{n-k} = x$ (where $k$
      is the number of digits in $d$) and the value round-trips.

\item \textbf{Formatting}: apply branch logic based on the relationship
      between $k$ (digit count) and $n$ (decimal exponent position) to
      produce the final string.  Key thresholds: $n \geq 1 \land n \leq
      21$ produces fixed notation; $n > 21$ or $n \leq 0$ with $|n| <
      certain bounds$ produces exponential notation.
\end{enumerate}

\subsection{IEEE 754 and Shortest Representation}

IEEE~754-2019~\cite{ieee754} binary64 has a 52-bit significand, 11-bit
exponent, and 1 sign bit.  The \emph{shortest representation} property
requires finding the decimal string with fewest digits that uniquely
identifies the original binary64 value among all representable values.
This is the core algorithmic challenge addressed by both \burgerdybvig
and \schubfach.
